%! Setting Up --- Learning from Data — Unit 8, Lesson 8.0 — Student Packet Cover Sheet
\documentclass[10pt]{article}
\usepackage{linalg-article}
\usepackage{linalg-boxes}
\usepackage{ltablex}
\keepXColumns

\begin{document}

% Full-bleed burgundy banner
\begin{tikzpicture}[remember picture, overlay]
  \fill[burgundy] (current page.north west)
    rectangle ([yshift=-0.9in]current page.north east);
\end{tikzpicture}
\vspace{-0.8in}

\begin{center}
  {\color{white}\textbf{\LARGE Linear Algebra}}\\[4pt]
  {\color{blushmid}\large\textbf{Unit 8 \quad Learning from Data}}\\[6pt]
  {\color{white}\textbf{\large Lesson 8.0 \quad Setting Up --- Learning from Data}}
\end{center}

\vspace{0.22in}
\namedateperiod
\vspace{0.18in}

\begin{learningtargetbox}
\small
\begin{enumerate}[leftmargin=1.5em, itemsep=5pt]
  \item \ldots{} explain that a \textbf{learning function} (a neural network) alternates a \textbf{linear step} (multiply by a weight matrix, add a bias) with the \textbf{nonlinear} ReLU $=\max(0,x)$.
  \item \ldots{} evaluate ReLU and compute one layer $\mathrm{ReLU}(A\mathbf{v}+\mathbf{b})$, and explain how stacking ReLUs builds \textbf{piecewise-linear} functions that bend to fit data.
  \item \ldots{} describe \textbf{learning} as adjusting the weights to shrink a \textbf{loss} (sum of squared errors) by \textbf{gradient descent}.
\end{enumerate}
\end{learningtargetbox}

\vspace{0.14in}

\begin{tocbox}
\small
\renewcommand{\arraystretch}{1.5}
\begin{tabularx}{\linewidth}{@{} c l X r @{}}
\rowcolor{blush}
\textbf{\#} & \textbf{Component} & \textbf{Description} & \textbf{Score} \\
\midrule
1 & Warm-Up        & Spiral review: matrix $\times$ vector, squared error, and the $\max(0,x)$ function & \blank{1.2cm} \\
2 & Guided Notes   & ReLU $=\max(0,x)$; one layer $\mathrm{ReLU}(A\mathbf{v}+\mathbf{b})$; piecewise-linear fits; loss \& gradient descent & \blank{1.2cm} \\
3 & Group Activity & Evaluating ReLU and one layer; building bent functions from ReLUs; comparing losses & \blank{1.2cm} \\
4 & Exit Ticket    & Quick check: one layer, a loss, and why the ReLU step is needed & \blank{1.2cm} \\
5 & Homework       & ReLU practice, layers with a bias, piecewise-linear graphs, comparing losses, and a look ahead & \blank{1.2cm} \\
\midrule
  & \multicolumn{2}{r}{\textbf{Total}} & \blank{1.2cm} \\
\end{tabularx}
\end{tocbox}

\vspace{0.10in}

\begin{remindbox}
\small A \textbf{neural network} is a \emph{function} that turns an input into a prediction, built from
tools you already know: each \textbf{layer} does a \emph{linear step} --- multiply by a \textbf{weight
matrix} and add a \textbf{bias} (Unit~1) --- then the one nonlinear step, the \textbf{ReLU}
$=\max(0,x)$, which keeps positives and zeroes negatives. That single fold, stacked and shifted, builds
\textbf{piecewise-linear} functions that bend to fit any data. The network \textbf{learns} by choosing its
weights to make a \textbf{loss} (the Unit~4 sum of squared errors) as small as possible, rolling downhill
by \textbf{gradient descent}. Modern AI is, at heart, the linear algebra of this course.
\end{remindbox}

\end{document}
